\section{Introduction}\label{sec:intro}

A \term{Latin square} of order $n$ is an $n\times n$ array filled with $n$ symbols so that each symbol appears exactly once in every row and every column. A \term{partial transversal} of size $m$ is a set of $m$ cells, no two in the same row or column, containing $m$ distinct symbols; a \term{transversal} is a partial transversal of size $n$. The cyclic Latin square $L(i,j)=i+j \bmod n$ has no transversal when $n$ is even (Claim~\ref{claim:cyclic}), so the following conjectures, due to Ryser~\cite{ryser1967} and, independently, Brualdi and Stein~\cite{stein1975}, would be best possible.

\begin{conjecture}[Ryser--Brualdi--Stein]\label{conj:rbs}
Every Latin square of order $n$ has a partial transversal of size $n-1$, and, if $n$ is odd, a transversal.
\end{conjecture}

It is convenient to translate this into the language of edge colourings. Given a Latin square $L$ of order $n$, colour the edge $ij$ of the complete bipartite graph $K_{n,n}$ with the symbol $L(i,j)$. The Latin property says exactly that this colouring is \term{proper}, that is, any two edges sharing a vertex receive different colours, and so every colour class is a perfect matching. A partial transversal of size $m$ is then a \term{rainbow matching} of size $m$: a matching whose edges all have different colours. Conjecture~\ref{conj:rbs} thus asserts that every proper $n$-edge-colouring of $K_{n,n}$ contains a rainbow matching with $n-1$ edges.

\topic{Previous work}

A simple greedy argument gives a rainbow matching with $\lceil n/2\rceil$ edges (see Theorem~\ref{thm:greedy} below). Koksma~\cite{koksma1969} improved this to $2n/3+O(1)$, and Woolbright~\cite{woolbright1978} and, independently, Brouwer, de~Vries and Wieringa obtained $n-\sqrt{n}$. Hatami and Shor~\cite{hatami2008} proved the bound $n-O(\log^2 n)$, and Keevash, Pokrovskiy, Sudakov and Yepremyan~\cite{kpsy2022} improved the error term to $O(\log n/\log\log n)$. Montgomery~\cite{montgomery2023} proved that every Latin square of sufficiently large order $n$ has a partial transversal of size $n-1$, settling the Brualdi--Stein part of Conjecture~\ref{conj:rbs} for large $n$.

\topic{These notes}

In Section~\ref{sec:greedy} we prove the greedy bound. The argument is short, but it contains the basic move behind almost every later improvement: an edge of an unused colour either extends the current rainbow matching or certifies that the matching is already large.
